% Bursting versus budding --- plan \S4.5.
% Sources: new_notes3 Part E \S22--\S25; new_notes Part II \S13; new_notes2
% \S9.  All displayed numerical values re-verified 2026-08-07.

\section{Bursting versus budding}
\label{p:sec:flooding}

Does bursting help a pathogen get established? The question belongs to the
stochastic early phase, where counts are small and extinction is a live
possibility \cite{pearson2011stochastic}, and here it can be answered
exactly:
both versions of the early infection are branching processes whose offspring
distributions are in closed form.

The mechanism at issue is worth naming before it is analysed. A bursting
cell delivers its whole output in one instant, so a founder that survives to
rupture presents the target population with many simultaneous infection
attempts rather than a trickle spread over its lifetime. Whether that
concentration helps or hurts is not obvious: many attempts at once buy
insurance against the loss of any one particle, but they are all forfeited
if the intracellular population dies before the burst. We call the
concentration \emph{flooding} throughout, and the question of this section
is when it pays.

\subsection{Offspring of one infected cell}

In the branching approximation each released particle independently either
infects a target cell, at rate $\gamma T$, or is cleared, at rate $c$, so
infects with probability
\begin{equation}
q=\frac{\gamma T}{\gamma T+c}.
\label{p:eq:q}
\end{equation}
The offspring of one infected cell is the burst size $\Kb$ thinned by that
probability. Summing \cref{p:eq:burstlaw}, the unconditional burst-size
\pgf{} is
$G_{\Kb}(z)=b+(\delta/\lambda)\,z/(a-z)$ --- including the mass $b$ of cells
that never burst --- and thinning replaces $z$ by $y:=1-q+qz$:
\begin{equation}
G_{\rm off}(z)=b+\frac{\delta}{\lambda}\,\frac{y}{a-y},
\qquad y=1-q+qz.
\label{p:eq:Goff}
\end{equation}
One checks $G_{\rm off}(1)=b+(\delta/\lambda)/(a-1)=b+B=1$. The mean and
variance of the offspring number are
\begin{equation}
m=G_{\rm off}'(1)=q\,V_\infty,
\qquad
\mathrm{Var}_{\rm burst}
=\frac{2q^2V_\infty}{a-1}+qV_\infty-q^2V_\infty^2,
\label{p:eq:offmoments}
\end{equation}
the variance following from $\mathrm{Var}(\Kb)$ via the limit of
\cref{p:eq:EW2} and the composition formula for random sums. When $\gamma
T\ll c$, $q\approx\gamma T/c$ and $m\approx R_0^{\mathrm{ODE}}$.

\subsection{Extinction probability in closed form}

The extinction probability of a branching process is the smallest root of
$G_{\rm off}(z)=z$ in $[0,1]$. With $y=1-q+qz$ the fixed-point equation is
quadratic, $z=1$ is always a root, and the other root gives, for $m>1$,
\begin{equation}
z_{\rm ext}^{\rm burst}
=\frac{a(1-q+qb)-1+q}{q}
=\frac{a-1}{q}+1-a(1-b).
\label{p:eq:zext}
\end{equation}

\subsection{Comparison with matched-mean budding}

The comparator is the budding cell of \cref{p:sec:classical}, whose output
is geometric with mean $p/d_{\Icell}$. Match it to the bursting cell's total
output, $p/d_{\Icell}=V_\infty$. That matching is what makes the comparison
fair, and by \cref{p:thm:R0inv} it costs nothing: at equal total output the
two mechanisms have exactly the same $R_0$, so everything that follows
compares two pathogens at one threshold, and any difference between them is
a difference in how the output is scheduled rather than in how much of it
there is. Thinning a geometric law leaves a geometric
law, so the budding offspring number is geometric with the same mean
$m=qV_\infty$, and its extinction probability is
\begin{equation}
z_{\rm ext}^{\rm bud}=\frac{1}{m}\qquad(m>1).
\label{p:eq:zextbud}
\end{equation}

Write $L:=a(1-b)=a-\mu/\lambda$, so that $V_\infty=L/(a-1)$.

\begin{theorem}[The flooding advantage]
\label{p:thm:flood}
For $m>1$,
\begin{equation}
z_{\rm ext}^{\rm burst}-z_{\rm ext}^{\rm bud}
=\bigl(L-1\bigr)\Bigl(\frac{1}{m}-1\Bigr).
\label{p:eq:flood}
\end{equation}
Hence bursting strictly lowers the establishment extinction probability,
relative to matched-mean budding, if and only if $L>1$, i.e.
\begin{equation}
\delta(\lambda+\mu)>\lambda\mu
\qquad\Longleftrightarrow\qquad
\boxed{\;\frac{1}{\delta}<\frac{1}{\lambda}+\frac{1}{\mu}\;}
\label{p:eq:floodcrit}
\end{equation}
(with the convention that the condition holds for all $\delta>0$ when
$\mu=0$). When $L=1$ the two extinction probabilities are identical for
every $q$.
\end{theorem}

\begin{proof}[Proof sketch]
The identity \cref{p:eq:flood} is algebra: substitute \cref{p:eq:zext},
$m=qV_\infty=qL/(a-1)$ and simplify. For the criterion, $L>1$ is $a(1-b)>1$,
i.e.\ $a>1+\mu/\lambda$; inserting the definition of $a$ and squaring (both
sides are positive, or the right-hand side is not) yields
$(\lambda+\mu+\delta)^2-4\lambda\mu>4\lambda^2(1+\mu/\lambda-\eta)^2$, which
reduces to $\delta(\lambda+\mu)>\lambda\mu$.
\end{proof}

\begin{figure}[tbp]
\centering
\includegraphics[width=0.64\textwidth]{N4b_4_L_landscape.pdf}
\caption{Where a pathogen sits relative to the flooding boundary. The
flooding parameter $L(\mu,\delta)$ at $\lambda=1$; the exact curve $L=1$,
equivalently $\delta_*=\lambda\mu/(\lambda+\mu)$, partitions the rate plane
into a region where bursting lowers extinction risk and one where it raises
it. The three marked points are the regimes of \cref{p:tab:trio},
$(\mu,\delta)=(0,0.1)$, $(0.5,1/3)$ and $(0.9,0.1)$, with $L=1.10$, $1.00$
and $0.42$.}
\label{p:fig:L-landscape}
\end{figure}

The identity is exact in each of the three regimes, and at both values of
$q$ tested on the boundary
(\cref{p:tab:flooding,p:fig:flooding-regimes}).

\begin{table}[htbp]
\centering
\renewcommand{\arraystretch}{1.1}
\begin{tabular}{@{}llccc@{}}
\toprule
$(\lambda,\mu,\delta)$ & $q$ & $z_{\rm ext}^{\rm burst}$ & $z_{\rm ext}^{\rm bud}$ & verdict \\
\midrule
$(1,0,0.1)$ & $0.2$ & $0.400$ & $0.455$ & bursting wins ($L=1.10$) \\
$(1,0.9,0.1)$ & $0.9$ & $0.935$ & $0.844$ & budding wins ($L=0.42$) \\
$(1,0.5,1/3)$ & $0.6$ & $0.833$ & $0.833$ & tie ($L=1$) \\
$(1,0.5,1/3)$ & $0.9$ & $0.556$ & $0.556$ & tie ($L=1$) \\
\bottomrule
\end{tabular}
\caption{Extinction probabilities from \cref{p:eq:zext,p:eq:zextbud}, and
the sign of their difference. The gap is $(L-1)(1/m-1)$ in every row.}
\label{p:tab:flooding}
\end{table}

\begin{figure}[tbp]
\centering
\includegraphics[width=0.94\textwidth]{F4b_2_flooding_regimes.pdf}
\caption{Extinction probabilities for bursting and matched-mean budding
across the three regimes of \cref{p:tab:trio}. Each pair of curves begins at
$m=1$; the signed gap follows \cref{p:eq:flood}. The $q$-windows differ
because $m=1$ sits at different $q$ in each regime.}
\label{p:fig:flooding-regimes}
\end{figure}

For $m\le1$ both processes die out almost surely and the comparison is
empty; \cref{p:eq:flood} changes sign there, an artefact of both formulae
running outside their domain $[0,1)$, not a reversal of the biology.

The flooding advantage is real, but it is not free. Bursting
concentrates output into one batch, lowering the offspring variance relative
to a budding cell's geometric output, but it carries the all-or-nothing
hazard that the intracellular population dies out first --- with probability
$b$ --- and the cell releases nothing. The criterion says the advantage
survives exactly when the burst timescale $1/\delta$ is shorter than the
harmonic-type combination $1/(1/\lambda+1/\mu)$ of the birth and death
timescales; when $\mu=0$ the hazard vanishes and bursting always wins.

\subsection{The boundary \texorpdfstring{$L=1$}{L=1}: a coincidence of laws}

The boundary case is more than a tie of extinction probabilities.

\begin{proposition}
\label{p:prop:L1}
The following are equivalent:
\begin{equation}
L=1
\;\Longleftrightarrow\;
a=1+\frac{\mu}{\lambda}
\;\Longleftrightarrow\;
\delta(\lambda+\mu)=\lambda\mu
\;\Longleftrightarrow\;
\frac{1}{\delta}=\frac{1}{\lambda}+\frac{1}{\mu}.
\label{p:eq:L1equiv}
\end{equation}
When they hold, the bursting offspring law \cref{p:eq:Goff} is
\emph{exactly geometric} with mean $m$ --- not merely equal in extinction
probability to the budding comparator. In particular the offspring variances
coincide, and
\begin{equation}
\mathrm{Var}_{\rm bud}-\mathrm{Var}_{\rm burst}
=\frac{2q^2V_\infty(L-1)}{a-1},
\label{p:eq:varorder}
\end{equation}
so that at matched mean $m$ the bursting offspring variance is smaller than
the budding one when $L>1$, equal when $L=1$, and larger when $L<1$.
\end{proposition}

\begin{proof}
$L=1$ is $a(1-b)=1$, i.e.\ $a-1=ab=\mu/\lambda$, the chain of equivalences
above. Under it, $b=\delta/\lambda$ (by \cref{p:eq:AB}), and
\cref{p:eq:Goff} collapses:
\[
G_{\rm off}(z)=b\Bigl(1+\frac{y}{a-y}\Bigr)=\frac{ab}{a-y}
=\frac{a-1}{a-y}
=\frac{(a-1)/a}{1-y/a},
\]
using $ab=a-1$. With $y=1-q+qz$,
\[
1-\frac{y}{a}=\frac{a-1+q}{a}\Bigl(1-\frac{q}{a-1+q}\,z\Bigr),
\]
so
\[
G_{\rm off}(z)=\frac{1-\varrho}{1-\varrho z},
\qquad
\varrho=\frac{q}{a-1+q}=\frac{m}{1+m},
\qquad m=\frac{q}{a-1}=qV_\infty,
\]
which is the \pgf{} of a geometric law on $\{0,1,2,\ldots\}$ with mean $m$
--- the same law as the matched budding comparator. The variance identity
\cref{p:eq:varorder} is the difference between the geometric variance
$m(1+m)$ and \cref{p:eq:offmoments}, with $V_\infty=L/(a-1)$.
\end{proof}

At the boundary, then, bursting and budding are the same branching process
in different clothes; away from it \cref{p:eq:varorder} quantifies how they
differ.

\subsection{The price: deterministic growth rate}
\label{p:sec:tradeoff}

Establishment is half the story; the other half is how fast an established
infection grows. Match the models a second time, at $R_0$ and at mean
productive lifetime: take the budding cell with
$d_{\Icell}=1/\mean{T_{\rm prod}}$ and $p=V_\infty d_{\Icell}$ --- exactly
$p_{\rm eff}(0)$, so the fairest classical match available --- and choose
$\gamma T$, or $c$, so both share $R_0=2$. Solving for the growth rates ---
$r_{\rm bud}$ from $(r+d_{\Icell})(r+c)=\gamma Tp$, $r_{\rm burst}$ from
\cref{p:eq:char} --- gives the same ordering in all three regimes of
\cref{p:tab:trio}:
\begin{equation}
r_{\rm bud}>r_{\rm burst}.
\label{p:eq:rorder}
\end{equation}
The values, with $c=1$, are \cref{p:tab:tradeoff,p:fig:growth-tradeoff}.

Whether \cref{p:eq:rorder} holds in general is not settled here, but it
reduces to a single inequality, in which the coupling constants do not
appear at all.

\begin{proposition}[The ordering is a Laplace-transform comparison]
\label{p:prop:rorder}
Match the comparator as above, so that
$\int_0^\infty g_{\rm bud}=\int_0^\infty g_{\rm burst}=V_\infty$. Then
$r_{\rm bud}>r_{\rm burst}$ for \emph{every} pair $(\gamma T,c)$ with
$R_0>1$ if and only if
\begin{equation}
\Lap{g}_{\rm burst}(r)\ \le\ \Lap{g}_{\rm bud}(r)=\frac{p}{r+d_{\Icell}}
\qquad\text{for all }r>0,
\label{p:eq:laplace-order}
\end{equation}
that is, if and only if the release age under bursting dominates the matched
exponential release age in the Laplace-transform order.
\end{proposition}

\begin{proof}
By \cref{p:eq:Atilde} each model's growth rate is the unique root of
$\phi(r):=\gamma T\Lap g(r)-(r+c)$; $\phi$ is strictly decreasing, and
$\phi(0)=\gamma TV_\infty-c=c(R_0-1)>0$, so the root is positive and unique.
If \cref{p:eq:laplace-order} holds then $\phi_{\rm burst}\le\phi_{\rm bud}$
pointwise, so the root of the former is no larger. Conversely, suppose
$\Lap g_{\rm burst}(r_0)>\Lap g_{\rm bud}(r_0)$ for some $r_0>0$. Choose any
$c>0$ and set $\gamma T=(r_0+c)/\Lap g_{\rm burst}(r_0)$, so that
$r_{\rm burst}=r_0$; then $\phi_{\rm bud}(r_0)<0$ and $r_{\rm bud}<r_0$.
Taking $c$ small makes $R_0=\gamma TV_\infty/c$ exceed one, so the ordering
fails for an admissible coupling.
\end{proof}

\noindent
Both sides of \cref{p:eq:laplace-order} agree at $r=0$, by the matching, so
its content is entirely in how the two release schedules are spread over
age. The inequality was checked at $1500$ triples drawn over
$(\mu,\delta)\in[10^{-3},10]\times[10^{-3},16]$ at $\lambda=1$, at $60$
values of $r$ across four decades, without a violation (test~N of
\cref{p:app:verification}). The \emph{usual} stochastic order does not hold
--- the burst release-age survival function dips below the matched
exponential --- so \cref{p:eq:laplace-order} is not a weaker way of saying
something stronger. It is the order the problem actually has.

\begin{table}[htbp]
\centering
\renewcommand{\arraystretch}{1.1}
\begin{tabular}{@{}lcccc@{}}
\toprule
$(\lambda,\mu,\delta)$ & $L$ & $d_{\Icell}$ & $r_{\rm bud}$ & $r_{\rm burst}$ \\
\midrule
$(1,0,0.1)$ & $1.10$ & $0.417$ & $0.250$ & $0.180$ \\
$(1,0.5,1/3)$ & $1.00$ & $0.910$ & $0.395$ & $0.294$ \\
$(1,0.9,0.1)$ & $0.42$ & $0.701$ & $0.343$ & $0.198$ \\
\bottomrule
\end{tabular}
\caption{Deterministic growth rates at matched $R_0=2$, matched total yield
and matched mean productive lifetime, with $c=1$. Budding invades faster on
both sides of the flooding boundary and on it.}
\label{p:tab:tradeoff}
\end{table}

\begin{figure}[tbp]
\centering
\includegraphics[width=0.66\textwidth]{F4b_3_growth_tradeoff.pdf}
\caption{The same comparison as \cref{p:tab:tradeoff}, drawn. Budding
invades faster than bursting in all three regimes of \cref{p:tab:trio},
which have $L=1.10$, $1.00$ and $0.42$: the ordering does not turn at the
flooding boundary, so it is not the same effect as
\cref{p:thm:flood}.}
\label{p:fig:growth-tradeoff}
\end{figure}

The mechanism is generation time in the exact sense of
\cref{p:sec:generation-kernel}: the budding cell starts releasing at once
while the bursting cell must first accumulate load, so budding's
generation-time density $\mathcal A/R_0$ sits earlier, and by
$\Lap{\mathcal A}(r)=1$ an earlier density gives a larger root. This is the
same delay that appeared as the eclipse-like lag of
\cref{p:fig:overlay-V,p:fig:growth-phase}, and it is the phage lysis-timing
problem in another notation: holding on longer makes more particles per cell
but takes longer to make them
\cite{wang2006lysis,wallinga2007generation}. That an infectious-period
distribution shifts the growth rate at fixed threshold is a general fact
about renewal epidemics rather than anything peculiar to bursting
\cite{lloyd2001realistic,champredon2015intrinsic}; what is particular here
is that the distribution is derived rather than assumed.
% NEEDS-REF: the optimal-lysis-timing literature beyond Wang (Abedon; Bull)

\begin{figure}[tbp]
\centering
\includegraphics[width=0.94\textwidth]{N4b_6_generation_times.pdf}
\caption{The generation-time structure behind the first row of
\cref{p:tab:tradeoff}. At $(\lambda,\mu,\delta)=(1,0,0.1)$, budding and
bursting are matched at $R_0=2$, $c=1$, total yield and mean productive
lifetime. The main curves are the normalised infection-generation laws
$\mathcal A(\alpha)/R_0$, including the common extracellular waiting-time
component; panel (b) isolates release age and shows why the budding law
leads. Bursting's mass sits in later generations: mean generation time
$3.398$ against $4.089$, which is what puts $r_{\rm bud}=0.250$ above
$r_{\rm burst}=0.180$.}
\label{p:fig:generation-times}
\end{figure}

Put together, the two comparisons say what bursting buys and what it costs.
At fixed total output per cell, bursting can lower the probability of early
extinction --- when $L>1$ --- while slowing deterministic invasion: $R_0$
stays fixed, and $z_{\rm ext}$ and $r$ move in opposite directions. Within
this comparison, then, a pathogen tuning $\delta$ faces a genuine trade-off
(\cref{p:fig:pareto}). For
$\mu>0$ budding beats bursting on extinction until $\delta$ crosses
$\delta_*=\lambda\mu/(\lambda+\mu)$, beyond which bursting takes over; for
$\mu=0$ bursting wins at every $\delta$, at the continuing cost of a slower
$r$.

\begin{figure}[tbp]
\centering
\includegraphics[width=0.62\textwidth]{N4b_5_pareto_extinction_growth.pdf}
\caption{Tuning $\delta$ traces a strict extinction--growth trade-off, drawn
parametrically as $\delta$ varies from $0.03$ to $0.8$ with $\lambda=1$,
$\mu=0$, $c=1$ and $R_0=2$. The coupling is reset along the curve to hold
$R_0=2$, and the budding comparator is matched in
lifetime yield $V_\infty$ and in mean productive lifetime through
$d_{\Icell}=1/\mean{T_{\rm prod}}$ and $p=V_\infty d_{\Icell}$. There is no
setting of $\delta$ that improves both coordinates: faster growth is bought
with a higher extinction probability.}
\label{p:fig:pareto}
\end{figure}

\begin{remark}[Scope of the comparison]
\label{p:rem:flood-scope}
Both results compare branching structures at fixed offspring mean. They say
nothing about nonlinear effects, where other advantages of flooding may
live; those are taken up in \cref{p:sec:limitations}. The evolutionary
framing, and where this result sits relative to the bursting--budding
spectrum, belongs to \ChPath{}, which inherits from this comparison the
awkward fact that one mechanism can be favoured on establishment and
disfavoured on invasion speed at the same time --- the nested-model problem
in its usual form \cite{gilchrist2006evolution}.
\end{remark}

\FloatBarrier
